\documentclass[aspectratio=169]{beamer}
\usetheme{Madrid}
\usecolortheme{seahorse}
\usepackage{graphicx}
\usepackage{booktabs}
\usepackage{xcolor}

\title{Laser Cutter Quick Start}
\subtitle{Trotec Q400 RF --- 60W CO\textsubscript{2}}
\date{\today}
\author{Good Studio --- Via Faentina 244/8, Firenze}

\begin{document}

\maketitle

% ---------------------------------------------------------------
\section{Background}

\begin{frame}{How a CO\textsubscript{2} Laser Cutter Works}
\begin{columns}
\begin{column}{0.55\textwidth}
\begin{itemize}
    \item RF-excited CO\textsubscript{2} tube generates a 10.6\,\textmu m infrared beam
    \item Beam is steered by mirrors and focused by a lens onto the work surface
    \item Focused spot ($\sim$0.1\,mm) vaporizes or melts material
    \item \textbf{Two modes:}
    \begin{itemize}
        \item \textbf{Engraving} --- beam rasters across the surface (like a printer)
        \item \textbf{Cutting} --- beam traces vector paths through the material
    \end{itemize}
    \item Air assist blows debris from the cut; exhaust extracts smoke
\end{itemize}
\end{column}
\begin{column}{0.4\textwidth}
\begin{block}{Key Parameters}
\begin{description}
    \item[Power] Beam intensity (\%)
    \item[Speed] Head travel rate (\%)
    \item[Frequency] Pulse rate (Hz)
    \item[Passes] Repeat count
\end{description}
\end{block}
\end{column}
\end{columns}
\end{frame}

% ---------------------------------------------------------------
\begin{frame}{What Can You Do With 60W?}
\begin{columns}
\begin{column}{0.5\textwidth}
\textbf{Cuts well:}
\begin{itemize}
    \item Acrylic --- up to $\sim$10\,mm
    \item MDF / HDF --- up to $\sim$6\,mm
    \item Plywood --- up to $\sim$6\,mm (depends on glue)
    \item Cardboard, paper, fabric
    \item Leather, felt, foam
    \item Thin wood veneers
\end{itemize}
\end{column}
\begin{column}{0.5\textwidth}
\textbf{Engraves well:}
\begin{itemize}
    \item Wood, MDF, cork
    \item Acrylic (frosted finish)
    \item Anodized aluminum
    \item Coated metals, glass (surface mark)
    \item Stone, slate, tile
\end{itemize}
\vspace{1em}
\alert{\textbf{Never cut:}}
\begin{itemize}
    \item PVC / vinyl (emits chlorine gas)
    \item Polycarbonate (burns, toxic)
    \item ABS (toxic fumes)
    \item Metals (reflects beam)
\end{itemize}
\end{column}
\end{columns}
\end{frame}

% ---------------------------------------------------------------
\section{Our Machine}

\begin{frame}{Trotec Q400 RF Overview}
\begin{columns}
\begin{column}{0.5\textwidth}
\begin{itemize}
    \item \textbf{Model:} Q400 RF (serial Q42-3066)
    \item \textbf{Tube:} 60W CeramiCore RF CO\textsubscript{2}
    \item \textbf{Bed:} 1000 $\times$ 610\,mm
    \item \textbf{Software:} Trotec Ruby (browser-based)
    \item \textbf{Location:} MiniPC\_103 controls the laser via USB
    \item \textbf{Network:} \texttt{MiniPC\_103} on shop Wi-Fi
\end{itemize}
\end{column}
\begin{column}{0.5\textwidth}
\begin{block}{Stock Sheets}
Our standard stock is 1000$\times$600\,mm, 6\,mm HDF.\\
Usable area: $\sim$950$\times$550\,mm (leave margin).
\end{block}
\vspace{0.5em}
\begin{block}{Tested Cut Settings (6\,mm HDF)}
\begin{tabular}{ll}
Power & 100\% \\
Speed & 0.3\% \\
Frequency & 1\,000\,Hz \\
Passes & 1 \\
Focus & Manual midplane \\
\end{tabular}
\end{block}
\end{column}
\end{columns}
\end{frame}

% ---------------------------------------------------------------
\section{Safety}

\begin{frame}{Safety Rules}
\begin{alertblock}{Exhaust Fan --- MANDATORY}
The vent fan \textbf{must} be running before you start \emph{any} job.
Laser cutting produces toxic smoke and particulates.
Running the laser without exhaust is dangerous to your health
and will contaminate the workspace.
\textbf{Verify airflow before every cut.}
\end{alertblock}

\begin{itemize}
    \item \textbf{Never leave a running job unattended.}
        Flare-ups can happen, especially with MDF and wood.
    \item Keep the lid closed during operation (interlock will stop the beam).
    \item Know where the fire extinguisher is.
    \item The beam is invisible (infrared) --- never bypass interlocks.
    \item Check material safety \emph{before} cutting. When in doubt, don't cut it.
\end{itemize}

\begin{block}{Shop Hours}
Laser use ends at \textbf{8:00\,PM}. The exhaust fan is loud ---
please be respectful of neighbors. Plan jobs to finish before cutoff.
\end{block}
\end{frame}

% ---------------------------------------------------------------
\section{Workflow}

\begin{frame}{Step 1: Prepare Your Design}
\begin{itemize}
    \item Create an \textbf{SVG} file (Inkscape, Illustrator, Fusion, etc.)
    \item Use \textbf{two colors} to separate operations:
\end{itemize}

\vspace{0.5em}
\begin{columns}
\begin{column}{0.45\textwidth}
\centering
\textcolor{red}{\rule{2cm}{3pt}}\\[4pt]
\textbf{Red} (\texttt{\#FF0000})\\
\textbf{= Cut}\\[4pt]
{\small Stroke color on vector paths.\\
Set stroke width very thin (0.01\,mm).\\
\texttt{fill="none"}}
\end{column}
\begin{column}{0.45\textwidth}
\centering
\rule{2cm}{3pt}\\[4pt]
\textbf{Black} (\texttt{\#000000})\\
\textbf{= Engrave}\\[4pt]
{\small Fill color for raster areas.\\
Text, logos, surface marks.\\
\texttt{stroke="none"}}
\end{column}
\end{columns}

\vspace{1em}
\begin{exampleblock}{Tip}
Set your document size in mm to match your stock.
Ruby imports at 1:1 scale when units are correct.
\end{exampleblock}
\end{frame}

% ---------------------------------------------------------------
\begin{frame}{Step 2: Import into Ruby}
\begin{enumerate}
    \item Log into the Windows NUC:\\
        {\small User: \texttt{Shamrock} / Password: \texttt{shamrock}}
    \item Open \textbf{Ruby} (desktop shortcut or \texttt{https://localhost:5001})
    \item Log into Ruby as admin:\\
        {\small Email: \texttt{r.j.walters@gmail.com} / Password: \texttt{Shamrock1}}\\
        {\small (Trotec requires a registered account --- just use Robb's)}
    \item Go to \textbf{Designs} and import your SVG
    \item Double-click the design to open the \textbf{Workbench}
    \item Assign the material (e.g.\ ``HDF 6mm'')
    \item Ruby maps red $\to$ Cut and black $\to$ Engrave automatically
\end{enumerate}

\vspace{0.5em}
\begin{alertblock}{Important}
If you change material settings, you must \textbf{delete the workbench}
and re-create it. Ruby caches parameters at workbench creation time.
\end{alertblock}
\end{frame}

% ---------------------------------------------------------------
\begin{frame}{Step 3: Focus \& Align}
\begin{enumerate}
    \item Place your stock on the bed
    \item \textbf{Focus the lens:}
    \begin{itemize}
        \item Use the physical gap tool (hangs on the machine)
        \item Place it between the lens assembly and the material surface
        \item Jog Z until the gap tool just fits --- this sets focus on the surface
    \end{itemize}
    \item \textbf{For thick cuts ($\geq$\,6\,mm):} manually raise the table $\sim$3\,mm
        after focusing to place the focal plane at the material midplane.\\
        {\small If the job also engraves, set the Engrave effect \textbf{ZOffset to +3\,mm}
        (table back down) so engraving focuses on the surface, not the midplane.}
    \item \textbf{Align the job:}
    \begin{itemize}
        \item Use Ruby's \emph{Prepare} screen to position the design on the bed
        \item The red pointer shows where the head is --- use it to check corners
    \end{itemize}
\end{enumerate}

\begin{exampleblock}{Focus matters}
Out-of-focus cuts are wider, slower to penetrate, and produce more char.
Always re-focus when changing materials or stock thickness.
\end{exampleblock}
\end{frame}

% ---------------------------------------------------------------
\begin{frame}{Step 4: Cut}
\begin{enumerate}
    \item Verify the \textbf{exhaust fan is running}
    \item Verify the \textbf{lid is closed}
    \item Press \textbf{Start} in Ruby (or the physical button on the machine)
    \item \textbf{Watch the job} --- stay nearby the entire time
    \item If you see sustained flame: press \textbf{Pause} or open the lid (kills beam)
    \item You can \textbf{pause and resume} at any time --- useful if you need to
        open the lid to adjust something, pull a finished part out while
        the rest is still cutting, or blow out a flare-up
\end{enumerate}

\vspace{1em}
\begin{block}{After the Job}
\begin{itemize}
    \item Let the exhaust run for 1--2 minutes to clear residual smoke
    \item Remove your parts and clean up debris / offcuts
    \item Wipe the bed if there's significant char residue
\end{itemize}
\end{block}
\end{frame}

% ---------------------------------------------------------------
\section{Tips \& Tricks}

\begin{frame}{Getting Good Results}
\begin{itemize}
    \item \textbf{Masking tape} on the surface reduces char marks on the face
    \item \textbf{Sand the edges} lightly after cutting to remove char
    \item \textbf{Kerf:} the beam removes $\sim$0.1--0.5\,mm of material.
        For press-fit joints, test with a kerf coupon first.
    \item \textbf{Taper:} with midplane focus, taper is minimal
        ($\sim$0.03\,mm difference top-to-bottom on 6\,mm HDF).
        Without midplane focus, expect $\sim$7\textdegree{} taper.
    \item \textbf{Inner geometry first:} Ruby cuts holes before outlines by default
        (good --- parts don't shift mid-cut)
    \item \textbf{Test first:} always run a small test piece before committing
        to a full sheet, especially with new materials.
\end{itemize}
\end{frame}

% ---------------------------------------------------------------
\begin{frame}{Project Ideas}
\begin{itemize}
    \item \textbf{Stack and castellate:} glue layers together and use interlocking
        butt joints to build thick, strong structures ---
        furniture, enclosures, jigs
    \item \textbf{Upholstery frames:} MDF makes excellent frames for
        upholstery work --- wrap with batting, stretch fabric over, staple
    \item \textbf{Composite wrapping:} tightly wrap cut MDF parts in packing tape
        for a smooth, sealed surface with added rigidity ---
        a quick ``poor man's laminate''
    \item \textbf{Fabric \& paper:} cut fabric, felt, cardstock, or sticker paper
        for textiles, packaging, stencils, and decals
    \item \textbf{Living hinges:} parallel kerf cuts in thin material
        let flat sheets bend into curves
    \item \textbf{Parametric design:} use code to generate SVGs ---
        finger joints, hex grids, test coupons, nesting layouts
\end{itemize}

\vspace{0.5em}
{\small This is a creative tool for the Good Studio community ---
if you have an idea, try it!}
\end{frame}

% ---------------------------------------------------------------
\section{Automation with Claude}

\begin{frame}{Using Claude Code with the Laser}
The laser can be controlled programmatically via SSH tunnels and the Ruby API.

\vspace{0.5em}
\begin{itemize}
    \item \textbf{Claude Code} can generate SVGs, configure materials,
        upload designs, and monitor jobs
    \item Useful for parametric designs, test sweeps, and batch operations
    \item The repo \texttt{trotec-tools} has the tunnel manager and Python client
\end{itemize}

\vspace{0.5em}
\begin{block}{Getting Started}
\begin{enumerate}
    \item Clone: \texttt{git clone github.com/project-shamrock/trotec-tools}
    \item Connect: \texttt{./tunnel.sh up}
    \item Open Claude Code in the repo directory
    \item Ask Claude to help with your design or cutting workflow
\end{enumerate}
\end{block}

\vspace{0.5em}
{\small Claude has context about the machine settings, API, material profiles,
and all the lessons learned from commissioning. It can read the docs and
pick up where previous sessions left off.}
\end{frame}

% ---------------------------------------------------------------
% ---------------------------------------------------------------
\section{Maintenance}

\begin{frame}{Keeping the Machine Happy}
\begin{itemize}
    \item \textbf{Lens cleaning:} the focus lens and mirrors accumulate smoke residue.
        Clean periodically with lens wipes and isopropyl alcohol.
        The optics are coated ZnSe --- handle gently, don't scratch.
        Look up Trotec lens cleaning on YouTube for technique.
    \item \textbf{Honeycomb tray:} MDF cutting deposits sticky resin on the tray.
        Remove it periodically and soak in warm water to dissolve the buildup.
    \item \textbf{Exhaust path:} check that the exhaust hose and fan intake
        aren't clogged with dust.
    \item \textbf{Replacement optics:} coated lenses are consumable.
        OEM replacements are expensive; compatible ZnSe lenses
        are available on AliExpress/Alibaba at lower cost.
\end{itemize}

\begin{block}{Trotec Support}
\textbf{Richard Welsing} --- Trotec Technical Support\\
\texttt{support@troteclaser.nl} --- Case\# 00532035 (or open new)
\end{block}
\end{frame}

% ---------------------------------------------------------------
\begin{frame}{Questions?}
\centering
\vspace{2em}
{\Large Happy cutting!}

\vspace{2em}
\begin{tabular}{rl}
\textbf{Machine} & Trotec Q400 RF, 60W CO\textsubscript{2} \\
\textbf{Software} & Trotec Ruby (browser) \\
\textbf{Materials} & HDF 6\,mm standard stock \\
\textbf{Shop hours} & Cutoff at 8\,PM \\
\textbf{Location} & Good Studio, Via Faentina 244/8, Firenze \\
\textbf{Support} & \texttt{support@troteclaser.nl} \\
\textbf{Repo} & \texttt{project-shamrock/trotec-tools} \\
\end{tabular}
\end{frame}

\end{document}
